% Part: lambda-calculus
% Chapter: syntax
% Section: eta

\documentclass[../../../include/open-logic-section]{subfiles}

\begin{document}

\olfileid{lam}{syn}{eta}
\olsection{تحويل~$\eta$}

ولحدود~$\lambda$ علاقة أخرى نبينها بالمثال الذي سبق في~\olref[fv]{sec}:
فالحد~$\lambd[x][(fx)]$ يأخذ وسيطًا ثم يطبق~$f$ عليه؛ فهو من حيث الدالة
نفس~$f$، إذ يختزل كل من~$\lambd[x][(fx)]N$ و$fN$ إلى~$fN$.
ولضبط هذا المعنى نستعمل اختزال~$\eta$ وتمديد~$\eta$.

\begin{defn}[تقلّص~$\eta$، $\eredone$]
  \ollabel{defn:beredone}
  علاقة \emph{تقلّص~$\eta$} ($\eredone$) أصغر العلاقات المتوافقة على الحدود
  التي تحقق
  \[
    \lambd[x][M x] \eredone M \text{ بشرط } x \notin \FV{M}
  \]
\end{defn}

\begin{defn}[اختزال~$\beta\eta$، $\bered$] \ollabel{defn:bered}
  علاقة \emph{اختزال~$\beta\eta$} ($\bered$) أصغر علاقة انعكاسية متعدية
  على الحدود تشتمل على~$\bredone$ و$\eredone$.
  فتعريفها يجمع قاعدتي الانعكاس والتعدي إلى القاعدتين:
  \begin{enumerate}
  \item من~$M \bredone N$ يلزم~$M \bered N$. \ollabel{defn:bered3}
  \item من~$M \eredone N$ يلزم~$M \bered N$. \ollabel{defn:bered4}
  \end{enumerate}
    
\end{defn}

\begin{defn}
  نضم إلى علاقة التكافؤ~$\equal$ قاعدة تحويل~$\eta$:
  \[
  \lambd[x][f x] \equal f \quad\text{إذا كان }x \notin \FV{f}
  \]
  ونسمي العلاقة الحاصلة~$\equal[\eta]$.
  % تصحيح مصدر (0366-I1): أُثبت شرط النضارة اللازم لتحويل إيتا.
\end{defn}

ومن شأن التكافؤ بحسب~$\eta$ أنه يتصل بامتدادية حدود لامبدا:

\begin{defn}[الامتدادية]
  نضم إلى علاقة التكافؤ~$\equal$ قاعدة~(\ext):
  \begin{center}
  من~$Mx \equal Nx$ يلزم~$M \equal N$، بشرط~$x \notin \FV{MN}$.
  \end{center}
  ونجعل~$\equal[\ext]$ رمزًا للعلاقة الحاصلة.
\end{defn}

ومؤدى هذه القاعدة، على وجه التقريب، أن الحدين إذا عدّا دالتين وكان
فعلهما في الوسيط الواحد سواء، وجب عدهما متساويين.

فلنثبت أن ما تضيفه قاعدة~$\eta$ هو الامتدادية وحدها، لا أكثر منها ولا أقل.

\begin{thm}
  $M \equal[\ext] N$ يكافئ~$M \equal[\eta] N$.
\end{thm}

\begin{proof}
  نبدأ بإثبات انغلاق~$\equal[\eta]$ تحت قاعدة الامتدادية؛ فإذا ثبت
  أن قاعدة~$ext$ لا تزيد شيئًا على~$\equal[\eta]$، لزم أن تحتوي~$\equal[\eta]$
  العلاقة~$\equal[\ext]$، وأن يلزم من~$M \equal[\ext] N$ تحقق~$M
  \equal[\eta] N$.

  ووجه انغلاق~$\equal[\eta]$ تحت~\ext{} أن كل نتيجة~$M \equal N$
  يراد اشتقاقها بقاعدة~\ext{} لها المقدمة~$Mx \equal[\eta] Nx$.
  ومنها نحصل على~$\lambd[x][Mx] \equal[\eta] \lambd[x][Nx]$
  بقاعدة من قواعد~$\equal$، ثم نطبق~$\eta$ على الطرفين فنبلغ~$M \equal[\eta] N$.

  وأما الاتجاه الآخر فنثبت فيه أن قاعدة~$\eta$ داخلة في~$\equal[\ext]$.
  فمتى أخذنا~$\lambd[x][Mx]$ و$M$ مع الشرط~$x \notin \FV{M}$،
  كان~$(\lambd[x][Mx])x \equal[\ext] Mx$؛
  فتنتج~$\lambd[x][Mx] \equal[\ext] M$ بقاعدة~\ext{}.
\end{proof}

\end{document}
